Theorem 4 (Connection to the Shapley Taylor interaction index). Given a subset of input variablesT ⊆N, thek-th order
Shapley Taylor interaction indexIShapley-Taylor(T ) can be represented as weighted sum of causal effects,i.e.,IShapley-Taylor(T ) =wT
if|T| <k ;IShapley-Taylor(T ) =∑
S⊆N\T
(|S|+k
k
)−1
wS∪T if|T| =k; andIShapley-Taylor(T ) = 0 if|T| >k .
• Proof: By the definition of the Shapley Taylor interaction index,
IShapley-Taylor(k)(T ) =



∆vT (x∅) if|T| <k
k
|N|
∑
S⊆N\T
1
(|N|−1
|S| )∆vT (xS) if|T| =k
0 if|T| >k
When|T| <k , by the definition of the Harsanyi dividend, we have
IShapley-Taylor(k)(T ) = ∆vT (x∅) =
∑
L⊆T
(−1)|T|−|L|·v(xL) =wT.
When|T| =k, we have
22
